% Generated draft: requires mathematical and source review.
\documentclass[11pt]{article}
\usepackage[margin=25mm]{geometry}
\usepackage{fontspec}
\setmainfont{Latin Modern Roman}
\usepackage{amsmath,amssymb,mathtools,unicode-math}
\setmathfont{Latin Modern Math}
\usepackage{enumitem,needspace,xurl,hyperref}
\setlength{\parindent}{0pt}
\setlength{\parskip}{5pt}
\setlength{\emergencystretch}{3em}
\newcommand{\N}{\mathbb{N}}
\newcommand{\Z}{\mathbb{Z}}
\newcommand{\Q}{\mathbb{Q}}
\newcommand{\R}{\mathbb{R}}
\newcommand{\C}{\mathbb{C}}
\newcommand{\F}{\mathbb{F}}
\newcommand{\A}{\mathbb{A}}
\newcommand{\PP}{\mathbb{P}}
\newcommand{\E}{\mathbb{E}}
\newcommand{\eps}{\varepsilon}
\newcommand{\dd}{\,\mathrm{d}}
\begin{document}
\section*{UnsolvedMath ID: 30004896}
\textbf{Source problem code:} catalogue entry
\subsection{Definitions and mathematical statement}
Write the potential in the chart denoted by \(T_i^\vee\) as a Laurent polynomial
\[
f=W|_{T_i^\vee}=\sum_{a\in A}c_a x^a,
\qquad x^a=\prod_{k=1}^{N}x_k^{a_k},
\]
where \(A\subset\Z^N\) is finite and \(c_a\in\C\setminus\{0\}\). The supplied excerpt also uses a chart \(T_\Sigma^\vee\) and asserts that all exponents of \(W|_{T_\Sigma^\vee}\) are nonpositive. It does not define the indexing data \(i,\Sigma\), their relationship, the potential \(W\), or the admissible family of charts.

Under coefficient-free minimum tropicalization, the Laurent polynomial determines the piecewise linear function
\[
f^{\mathrm{trop}}(u)=\min_{a\in A}\langle a,u\rangle,
\qquad u\in\R^N.
\]
This specifies the convention in which nonzero coefficients contribute no tropical constant. The excerpt describes multiplication becoming addition and addition becoming minimum, but does not supply a coefficient valuation or a complete definition of the resulting inequality system. In particular, its truncated formula does not reliably specify the tropical domain.

To express redundancy without supplying missing inequality conventions, let \(D\) denote the intended ambient domain and let \(I_a(u)\) denote the inequality associated with the term \(c_ax^a\). Set
\[
P=\{u\in D:I_b(u)\text{ holds for every }b\in A\}.
\]
The inequality \(I_a\) is redundant precisely when deleting it leaves \(P\) unchanged, equivalently when
\[
\forall u\in D,\qquad
\left(\bigwedge_{b\in A\setminus\{a\}}I_b(u)\right)
\Longrightarrow I_a(u).
\]
Any additional fixed constraints must be incorporated into \(D\).

\textbf{Recorded conjecture, with unresolved qualifications.} The clean statement and the extracted Conjecture 4 ask whether, for every admissible potential and chart and every term \(c_ax^a\),
\[
I_a\text{ is redundant}\quad\Longleftrightarrow\quad c_a\leq 1.
\]
This is a schematic formulation: the admissible data, the ambient domain, the individual inequalities, and an ordered coefficient domain remain unspecified. The displayed comparison requires real coefficients or another explicitly ordered coefficient domain; membership in a complex Laurent polynomial ring alone does not define it.

\textbf{Conflicting supplied formulation.} The background and literature assessment instead describe the criterion
\[
I_a\text{ is redundant}\quad\Longleftrightarrow\quad c_a>1.
\]
For real coefficients, these are incompatible proposed characterizations. The supplied evidence does not resolve which comparison is intended, so neither is selected as a corrected conjecture.

\textbf{Reported background qualification.} The excerpt labels as Theorem 3, attributed there to reference [3], the assertion that if every variable exponent in \(W|_{T_i^\vee}\) belongs to \(\{-1,0\}\), then the inequalities of \(C_{\Sigma_i}\), and hence those of \(S_i\), are nonredundant. The objects \(C_{\Sigma_i}\), \(S_i\), the full hypotheses, and the identity of reference [3] are not supplied. This background assertion is retained as reported context, not as an independently verified theorem or a resolution of the conflicting coefficient criteria.
\subsection{Short English statement}
The question asks whether the coefficient of a monomial in a chart expression of the potential exactly determines whether its tropical inequality can be deleted without changing the feasible set. The recorded conjecture uses coefficient at most \(1\), while the supplied literature assessment uses coefficient greater than \(1\). The intended comparison, coefficient assumptions, chart data, and inequality system require clarification before a definitive mathematical statement can be accepted.
\subsection{Sources}
\begin{itemize}
\item Supplied local UnsolvedMath registry record, permanent ID \(30004896\), source problem code \texttt{catalogue entry}, titled ``Coefficient Criterion for Redundant Tropical String Inequalities.'' The record gives proposed year 2021 and does not identify a proposer.
\item Supplied report citation: \emph{Enveloping Algebras and Geometric Representation Theory} (2022), Oberwolfach Reports, pp.\ 2827--2891, DOI \texttt{10.4171/owr/2021/52}. \url{https://doi.org/10.4171/owr/2021/52}. The supplied excerpt labels the target Conjecture 4 and refers to a preceding Theorem 3 with an unresolved reference [3].
\item Supplied literature source: Koshevoy--Schumann, \emph{Journal of Algebraic Combinatorics}. \url{https://doi.org/10.1007/s10801-022-01144-z}. No article title, precise result location, or full bibliography is supplied.
\item Supplied literature source: ``Oberwolfach Report 52, string cones.'' \url{https://doi.org/10.4171/owr/2021/52}.
\item Literature assessment dated exactly \texttt{2026-08-22}: the supplied curation labels the problem open and states that the coefficient-greater-than-one criterion remains conjectural, with examples and sufficient criteria but no general proof or counterexample in the primary paper. This is supplied dated curation, not independently verified literature assessment; its coefficient comparison conflicts with the recorded conjecture.
\item The registry labels the statement \texttt{corrected\_verified} and reports that the official report was checked directly. That verification claim is supplied provenance and has not been independently verified in this draft.
\end{itemize}
\end{document}
